\chapter{算法模板与工程清单}

本章把前文的概念整理为可以直接改写为实验代码的流程。伪代码不是某个框架的完整实现，它的作用是暴露数据、优化和评估之间的边界。

\section{监督学习训练模板}
\begin{algorithm}[htbp]
\caption{带验证集和早停的监督学习}
\begin{algorithmic}[1]
\Require 训练集$\mathcal{D}_{tr}$，验证集$\mathcal{D}_{va}$，模型$f_\theta$
\State 初始化$\theta$，设置最佳验证损失$L^*=+\infty$
\For{$e=1,\ldots,E$}
 \For{每个小批次$B\subset\mathcal{D}_{tr}$}
  \State $L_B\gets |B|^{-1}\sum_{(x,y)\in B}\ell(f_\theta(x),y)$
  \State $\theta\gets\operatorname{OptimizerStep}(\theta,\nabla_\theta L_B)$
 \EndFor
 \State 计算验证损失$L_e$，若$L_e<L^*$则保存检查点
 \State 若连续$p$轮没有改善则停止
\EndFor
\Ensure 返回验证集最优检查点，而非最后一轮参数
\end{algorithmic}
\end{algorithm}

早停的关键不是“训练少几轮”，而是用未参与参数更新的数据选择模型。若反复根据测试集结果调整超参数，测试集就不再是无偏的最终评估。

\section{时间序列滚动评估}
给定时间序列，可使用扩展窗口：第$k$次实验用$[1,t_k]$训练，在$(t_k,t_{k+1}]$测试。每个窗口的模型可以重新训练，也可以模拟在线更新。必须记录窗口之间是否允许使用新标签，以及标准化统计量是否随窗口更新。

\section{联邦学习模板}
\begin{algorithm}[htbp]
\caption{联邦平均的抽象流程}
\begin{algorithmic}[1]
\Require 初始模型$\theta^0$，客户端集合$\mathcal{K}$，通信轮数$T$
\For{$t=0,\ldots,T-1$}
 \State 服务器抽样客户端集合$S_t$
 \For{每个$k\in S_t$并行执行}
  \State 客户端从$\theta^t$开始本地训练$E$步，得到$\theta_k^{t+1}$
  \State 上传更新或安全聚合所需的掩码更新
 \EndFor
 \State 服务器按样本量聚合$\theta^{t+1}=\sum_{k\in S_t}p_k\theta_k^{t+1}$
\EndFor
\end{algorithmic}
\end{algorithm}

工程实现必须处理客户端版本不一致、上传超时、异常更新和聚合权重归一化。若只在理想同步环境中测试，不能声称系统满足真实联邦约束。

\section{实验配置的分层}
建议把配置分成四层。数据层规定路径、字段、时间范围和划分；模型层规定结构、宽度、深度和预训练权重；优化层规定批大小、学习率、调度和随机种子；评估层规定指标、阈值、分组字段和输出目录。分层配置能减少“改了一个参数却不知道影响了什么”的问题。

\section{数据版本与特征血缘}
每个训练特征都应能追溯到原始字段、变换代码、窗口定义和时间可得性。特征血缘对于发现数据泄漏和复现历史模型尤其重要。建议保存训练样本的唯一标识和特征生成版本，而不是只保存最终矩阵。

\section{监控指标}
上线监控至少包括输入缺失率、数值范围、类别新值、推理延迟、资源使用、置信度分布和人工反馈。标签到达后补充分层性能、校准和错误成本。监控系统应区分瞬时异常、持续漂移和数据管道中断，避免把所有问题都归因于模型。

\section{模型回滚与影子部署}
影子部署让新模型接收真实流量但不影响决策，可比较延迟、输入覆盖和预测分布。灰度发布按设备或流量比例逐步扩大。回滚条件应在上线前定义，例如延迟超阈值、关键类别漏报上升或输入字段大面积缺失。回滚后保留事件日志，便于定位是模型、数据还是服务问题。

\section{安全与权限}
训练数据、模型检查点、检索知识库和日志具有不同敏感级别。模型服务应限制管理接口，校验输入大小，隔离用户提供的文本与系统指令，并记录高风险操作。提示词过滤不能替代权限控制；能否执行动作应由独立的策略层决定。

\section{论文结果表的最低要求}
结果表应包含任务、数据划分、输入模态、模型参数量、训练预算、均值和波动范围。若不同方法使用不同外部数据、不同预处理或不同测试集，应在表注中说明。显著数字不应掩盖样本量很小或测试集被反复使用的问题。

\section{提交前检查清单}
\begin{enumerate}[label=\arabic*.]
 \item 数据划分文件是否固定并可追溯？
 \item 训练和推理是否使用完全一致的预处理？
 \item 所有基线是否获得合理的超参数预算？
 \item 结果是否来自多次运行或给出统计不确定性？
 \item 是否进行了设备、时间、缺失和噪声分层评估？
 \item 是否保存了最佳检查点和完整配置？
 \item 失败案例是否被单独分析，而不是删除？
 \item 部署时的延迟、资源、权限和回滚策略是否经过测试？
\end{enumerate}

\section{进一步练习}
\begin{enumerate}[label=\arabic*.]
 \item 将监督学习伪代码改写为带梯度累积的版本，并说明等效批大小。
 \item 为滚动预测设计一个防止未来信息进入标准化统计量的实现方案。
 \item 为联邦学习加入客户端掉线和最小参与数条件。
 \item 设计一个影子部署报告，包含输入漂移、延迟和预测一致性。
\end{enumerate}

\section{数据管道伪代码}
\begin{algorithm}[htbp]
\caption{带时间边界的数据准备}
\begin{algorithmic}[1]
\Require 原始事件流$\mathcal{E}$，预测时刻集合$\mathcal{T}$，历史窗口$w$
\For{每个预测时刻$t\in\mathcal{T}$}
 \State 读取时间区间$[t-w,t]$内可用的观测
 \State 删除在$t$之后才产生的字段和事后修订字段
 \State 按训练阶段保存的统计量进行变换
 \State 根据预先冻结的标签协议生成$y_t$
 \State 输出样本标识、设备标识、时间戳、特征和标签
\EndFor
\State 按设备、批次或时间执行划分，并保存划分文件
\Ensure 数据处理日志、异常记录和特征版本
\end{algorithmic}
\end{algorithm}

该模板强调“可用信息”的边界。工程代码中最容易被忽视的不是模型类，而是时间过滤和统计量保存。每一个特征都应能回答：它在预测时刻是否已经产生？它是否经过未来信息计算？它是否在训练和推理中使用同一单位？

\section{主动学习循环}
当标注成本很高时，可从未标注池中选择最值得标注的样本。选择函数可以基于预测熵、模型间分歧、代表性或业务风险。一个基本循环为：
\begin{enumerate}[label=\arabic*.]
 \item 用当前标注集训练模型；
 \item 在未标注池上计算不确定性和覆盖度；
 \item 结合业务约束选择一批样本；
 \item 由专家标注并记录标注时间和意见冲突；
 \item 更新标注集，重新训练并在固定测试集上评估。
\end{enumerate}
主动学习不应根据测试集错误直接挑选样本，否则测试集会参与训练过程。还需比较随机采样、分层采样和主动采样的标注效率，回答每增加一单位标注成本带来多少有效收益。

\section{模型压缩与边缘部署}
工业现场可能只有有限显存和功耗预算。剪枝删除冗余连接，量化使用低精度表示，知识蒸馏让小模型拟合教师模型的软输出。压缩后应同时测量精度、延迟、峰值内存、功耗和启动时间；服务器上的 FLOPs 不能直接代表边缘设备上的实际延迟。

若量化导致某些罕见故障类别性能明显下降，整体平均指标可能仍然看不出来。因此压缩验收必须分层报告，并对安全关键类别设置硬约束。模型压缩也可能改变概率校准，阈值需要重新验证，不能直接沿用压缩前阈值。

\section{检索增强系统评估模板}
对于问答系统，至少拆成四个层次：检索召回率、证据排序、答案事实性和引用完整性。可以构造包含标准答案、有效证据、不可回答问题和冲突文档的评估集。对不可回答问题，系统应在没有足够证据时拒答，而不是凭语言流畅度给出看似合理的答案。

\section{故障排查顺序}
当实验结果异常时，建议按以下顺序排查：首先确认数据行数、标签分布和划分哈希；其次检查一个小批次能否过拟合；再次检查损失、梯度和学习率；然后比较无增强、无预训练和无复杂模块的基线；最后才分析网络结构。这个顺序能把数据错误、实现错误和建模错误分开。

\begin{table}[htbp]
\centering
\caption{常见现象与优先检查项}
\begin{tabular}{p{0.27\textwidth}p{0.59\textwidth}}
\toprule
现象 & 优先检查\\
训练损失不下降 & 标签类型、输出激活、学习率、梯度是否为零\\
训练很好验证很差 & 划分、重复样本、设备泄漏、过拟合和标签噪声\\
结果每次差异很大 & 随机种子、批次大小、类别采样、数据顺序和归一化\\
上线后输入异常 & 字段版本、单位、时区、缺失编码和服务序列化\\
告警数量突然变化 & 输入漂移、阈值、标签反馈和规则版本\\
\bottomrule
\end{tabular}
\end{table}

\section{研究计划模板}
一个四周的复现实验可以这样安排：第一周固定数据和指标，复现最简单基线；第二周复现主方法并核对训练曲线；第三周完成关键消融和跨设备测试；第四周整理误差分析、资源开销和限制。每周都应产生可检查的中间产物，不要把所有风险留到最后一晚的总实验。

\section{结果解读与研究答辩}
研究汇报不应只展示一张最高分表格。首先说明任务、数据划分和评价指标；然后说明基线与当前方法究竟保持了哪些条件一致；最后解释性能差异可能来自结构、数据、训练预算还是随机性。一个可信的结论应与实验设计一一对应，例如“在按设备划分的测试中，加入频域特征后事件级召回提高”，而不是泛化为“该方法具有更强泛化能力”。

回答质疑时，可以把问题分为三类。对于事实问题，直接给出数据、代码或日志位置；对于方法问题，说明目标函数和设计假设；对于外推问题，明确当前实验没有验证的范围。承认限制不会削弱研究，反而能帮助读者判断结论适用的边界。

结果表中的小数位数也应与实验精度匹配。如果不同随机种子的标准差为 $0.8$ 个百分点，就不应把均值写成看似精确到四位小数的结论。对于多个数据集和多个指标，应预先确定主要指标，避免在结果出来后只挑选最有利的一个。若进行了多次测试集选择，应如实说明并重新建立独立的最终评估。

一个完整的答辩结论可以采用如下结构：本工作解决了什么实际约束；采用了什么最小必要方法；在什么严格划分下获得了什么变化；哪些消融支持机制解释；在哪些条件下失败；下一步需要什么数据或实验。这个结构能把“模型做了什么”与“证据支持什么”分开，也是阅读论文和撰写论文时最值得反复练习的表达方式。